\chapter{Използвани технологии}

\section{View технологии}

\subsection{Blazor}

Blazor е избран като уеб технология в проекта, тъй като позволява изграждане на интерактивен потребителски интерфейс изцяло в .NET среда, без необходимост от отделен приложен слой на JavaScript. За конкретната разработка това се оказва особено удобно, защото споделените ViewModel-и могат да бъдат инжектирани директно в компонентите и използвани като основен източник на състоянието на интерфейса. Освен това Blazor е много подходящ за изграждане на преизползваеми елементи на интерфейса, каквито в проекта са панелът за филтриране и компонентът за визуализация на списъци.

Blazor е полезен не само като технология за потребителски интерфейс, а и като добра среда за изграждане на автоматизирано генериран потребителски интерфейс върху общи приложни абстракции. Компонентният модел на Blazor улеснява изграждането на повторно използваеми визуални елементи, които работят върху метаданни и ViewModel-и, а не върху конкретни класове от модела. Това го прави подходящ избор за филтри, таблици и детайлни изгледи, които се повтарят в различни части на приложението.

\subsection{AvaloniaUI}

Avalonia е избрана като настолна технология, защото предоставя XAML, модел на обвързване и команди, близък до класически MVVM инструменти като WPF, но без да бъде ограничена само до Windows. Това е важен практически аргумент, тъй като проектът се разработва и тества под Linux. В тази среда Avalonia позволява реална настолна реализация със сравнително тънък code-behind и естествено обвързване със споделен ViewModel слой.

В контекста на настоящия проект Avalonia е използвана като пълноценен настолен клиент, а не само като минимална визуална обвивка. Настолният клиент реално поддържа обзорни метрики, преглед на видео библиотеката, качване на видеа, стартиране на анализи, разглеждане на историята от изпълнения и визуализация на резултатите. По този начин се демонстрира, че споделените приложни ViewModel-и могат да обслужват втори интерфейс с реална функционалност.

\section{Системи за управление на бази данни}

\subsection{SQLite}

SQLite е подходяща за локална разработка и за демонстрационни сценарии, тъй като не изисква отделен сървър и работи чрез един файл. Това я прави удобна за бързо стартиране на приложението, за преносима демонстрационна конфигурация и за ситуации, в които е важно системата да може да бъде показана без допълнителна инсталация на сървър за база данни. В рамките на проекта SQLite е естествен избор за лек режим на работа и за лесна повторяемост при локално тестване.

Същевременно SQLite има и ограничения. Тя не отразява напълно типична експлоатационна среда и не е най-подходящият избор за по-сериозен многопотребителски сценарий. Именно поради това в проекта е включена и PostgreSQL, така че слоят за съхранение да не бъде обвързан само с една удобна за разработка технология.

\subsection{PostgreSQL}

PostgreSQL е избрана като втора СУБД, защото е значително по-близка до реална сървърна среда. Тя предоставя по-богата екосистема, по-реалистичен модел на внедряване и добра интеграция с EF Core чрез Npgsql. В контекста на проекта PostgreSQL показва, че една и съща приложна и бизнес логика може да работи върху различни доставчици без дублиране на код.

Най-същественото тук е, че проектът не поддържа две бази чрез два отделни service слоя. Вместо това същият \texttt{AppDbContext}, същите инфраструктурни услуги и същите приложни ViewModel-и работят и с SQLite, и с PostgreSQL. Разликата се свежда до конфигурационен избор на доставчик, а не до паралелни имплементации със същата логика.

\section{Други използвани технологии}

\subsection{Entity Framework Core}

Entity Framework Core е основата на слоя за съхранение. Чрез него проектът получава общ модел на данните, общ механизъм за заявки и възможност за превключване между SQLite и PostgreSQL. За конкретната разработка EF Core е ценен и с това, че позволява приложната логика да остане сравнително чиста и независима от конкретния доставчик на база данни, докато инфраструктурният слой поема конфигурацията и реалната интеграция с базата.

\subsection{YoloDotNet, ONNX Runtime и ffmpeg}

Важна особеност на проекта е, че частта за видео анализ не е оставена на ниво симулация. Вместо фиктивни записи за разпознати обекти е реализирана работеща последователност с \texttt{YoloDotNet}, ONNX модел и \texttt{ffmpeg}. \texttt{ffmpeg} се използва за извличане на подбрани кадри и за рендериране на анотирано изходно видео, а YoloDotNet изпълнява разпознаване на обекти върху тези кадри. Именно тази комбинация превръща приложението от обикновен административен интерфейс в система с реална аналитична функционалност.

Официалният проект на YoloDotNet го представя като модулна .NET библиотека за YOLO инференция, която стъпва върху ONNX Runtime и е насочена към предвидимо и производително изпълнение без Python среда за изпълнение и без тежки външни библиотеки за компютърно зрение \cite{yolodotnet}. Тази идея съответства много добре на целите на настоящия проект, защото \projectname{} също се стреми към изцяло .NET реализация, която да бъде лесна за демонстрация и достатъчно компактна за учебна среда.

Разпознаването на обекти представлява задача от компютърното зрение, при която системата трябва едновременно да определи какъв обект присъства в изображението и къде точно се намира той. За разлика от обикновената класификация, тук не е достатъчно да се каже, че в кадъра има човек или автомобил. Необходимо е и локализиране на всеки обект чрез ограничителна рамка, както и оценка на увереността на модела. Именно този тип резултат се съхранява в \projectname{} под формата на етикет, увереност и координати на ограничителната рамка.

\optionalfigure[0.78\textwidth]{figures/branding/object_detection_example.png}{Пример за резултат от разпознаване на обекти с bounding boxes}{fig:object-detection-example}

В основата на много modern computer vision модели стоят конволюционните невронни мрежи. Тяхната идея е, че чрез конволюционни филтри могат да се извличат локални визуални признаци като контури, текстури и по-сложни пространствени структури. С нарастване на дълбочината на мрежата тези признаци стават все по-абстрактни и подходящи за разпознаване на цели обекти. Макар съвременните YOLO архитектури да съдържат и допълнителни оптимизации, връзката им с CNN подхода остава фундаментална, защото те също използват learnable feature extraction върху изображението, за да достигнат до финално предсказване.

\optionalfigure[0.72\textwidth]{figures/branding/CNN.png}{Схематично представяне на конволюционна невронна мрежа}{fig:cnn-overview}

YOLO, съкратено от \textit{You Only Look Once}, е семейство от едностъпкови модели за разпознаване на обекти. Основната идея е, че вместо да се правят отделни етапи за генериране на предложения за региони и последваща класификация, моделът обработва изображението в един общ проход за инференция и директно предсказва класове и ограничителни рамки. Това го прави особено подходящ за приложения, в които е важен балансът между точност и скорост, например обработка на видео и почти реално време.

\optionalfigure[0.82\textwidth]{figures/branding/Yolo_Algorithm.png}{Схема на общата идея зад YOLO алгоритмите}{fig:yolo-algorithm}

В проекта е избрана серията YOLOv8. Причината не е единствено популярността на модела, а добрият практически баланс между зрялост на екосистемата, стабилна документация, лесен export към ONNX и широка поддръжка в инструменти като Ultralytics и YoloDotNet. Официалната екосистема на Ultralytics представя YOLOv8 като state-of-the-art поколение, създадено за задачи като detect, segment, classify и pose, с няколко размерни варианта, които позволяват избор според конкретния хардуер и натоварване \cite{ultralyticsyolov8, ultralyticsmodels}.

По-конкретно в \projectname{} се използва \texttt{yolov8n.onnx}, тоест nano вариантът на YOLOv8. Този избор е напълно логичен за учебен проект, който трябва да работи надеждно на стандартна машина и в настолно приложение с CPU доставчик за изпълнение. Според публикуваните от Ultralytics сравнения nano моделът е най-малкият вариант по брой параметри и изчислителна сложност, като същевременно остава практически полезен за реални задачи по детекция \cite{ultralyticsmodels}. С други думи, избраният модел не максимизира абсолютната точност, а търси добър инженерeн компромис между време за инференция, използвана памет и достатъчно добри резултати за демонстрация върху видео.

Този избор на технологии носи и практическа стойност за документацията. Проектът показва не само UI и работа с база данни, а и интеграция с външен инструмент и с библиотека за машинно зрение. Следователно той обхваща по-широк спектър от инженерни задачи, отколкото стандартен учебен пример.
